% Generated draft: requires mathematical and source review.
\documentclass[11pt]{article}
\usepackage[margin=25mm]{geometry}
\usepackage{fontspec}
\setmainfont{Latin Modern Roman}
\usepackage{amsmath,amssymb,mathtools,unicode-math}
\setmathfont{Latin Modern Math}
\usepackage{enumitem,needspace,xurl,hyperref}
\setlength{\parindent}{0pt}
\setlength{\parskip}{5pt}
\setlength{\emergencystretch}{3em}
\newcommand{\N}{\mathbb{N}}
\newcommand{\Z}{\mathbb{Z}}
\newcommand{\Q}{\mathbb{Q}}
\newcommand{\R}{\mathbb{R}}
\newcommand{\C}{\mathbb{C}}
\newcommand{\F}{\mathbb{F}}
\newcommand{\A}{\mathbb{A}}
\newcommand{\PP}{\mathbb{P}}
\newcommand{\E}{\mathbb{E}}
\newcommand{\eps}{\varepsilon}
\newcommand{\dd}{\,\mathrm{d}}
\begin{document}
\section*{UnsolvedMath ID: 30000951}
\textbf{Source problem code:} catalogue entry
\subsection{Definitions and mathematical statement}
The source poses three related growth questions.
\begin{enumerate}
\item In the usual finite undirected graph setting with at least two vertices, let \(M\) be the adjacency matrix and order its eigenvalues algebraically as \(\lambda_1\geq\lambda_2\geq\cdots\). Obtain lower bounds for the adjacency spectral gap \(\lambda_1-\lambda_2\). The original question explicitly asks for lower bounds; no graph class, degree, size parameter, normalization, or proposed bound is specified.
\item For a finite subset \(A\) of an abelian group written additively, define
\[
A+A=\{a+b:a,b\in A\}.
\]
How large is \(|A+A|\) compared with \(|A|\), and what bounds or structural explanations are available?
\item For a finite subset \(A\) of an arbitrary group written multiplicatively, define
\[
AAA=\{abc:a,b,c\in A\}.
\]
How large is \(|AAA|\) compared with \(|A|\), and why?
\end{enumerate}
For the two set-growth questions the nonempty case is the substantive one; no assumption that \(A\) generates its ambient group, is symmetric, or avoids proper subgroups is supplied. These are broad requests, not a single specified expansion theorem. In particular, disconnected graphs and finite subgroups are not excluded by the literal wording, so a uniform positive gap or a uniform strictly superlinear growth conclusion would require additional hypotheses. The record does not select such hypotheses, a target exponent, or an exceptional structured class.
\subsection{Short English statement}
Seek lower bounds for adjacency spectral gaps and understand sumset and triple-product growth relative to set size. The graph and group hypotheses needed for a nontrivial general bound are not specified.
\subsection{Sources}
\begin{itemize}
\item Local UnsolvedMath registry, permanent ID \(30000951\), source code catalogue entry, ``Spectral and Product Growth in Graphs and Groups.'' No individual proposer is supplied.
\item \emph{Groups and Geometries} (2009), Oberwolfach Reports, DOI 10.4171/owr/2008/20, pp. 1031--1078, also listed as ``Oberwolfach Report 20/2008'': \url{https://doi.org/10.4171/owr/2008/20}.
\item Supplied reference, ``Helfgott, growth in SL3'': \url{https://arxiv.org/abs/0807.2027}.
\item Supplied reference, ``Helfgott, growth in SL2'': \url{https://doi.org/10.4007/annals.2008.167.601}.
\item Supplied literature triage, checked 2026-08-21: status ``partially solved.'' It credits Helfgott and others with major product-growth and spectral-gap results for important finite groups including \(\operatorname{SL}_2\) and \(\operatorname{SL}_3\), while describing the broad questions as open. Precise hypotheses are not supplied. This dated assessment derives from supplied OWR curation; time-sensitive status and cited claims require independent verification.
\end{itemize}
\end{document}
